% Unit 030 terjemahan Bahasa Indonesia dari chapter2.tex baris 1756--2132.
% Sumber: Methods of Algebra, Volume 2, commit
% 9a5803ff77dd3257484cb177f851a73770a59dd3 (CC BY 4.0).
% stable-unit-id: o014.aljabr2.chapter2.grothendieck-categories
% Cakupan menutup Bagian 2.10 dan seluruh latihan Bab 2.

% segment-id: o014.li2.chapter2.grothendieck-categories.h001
\section{Kategori Grothendieck}\label{sec:Grothendieck-cat}
\hypertarget{o014.aljabr2.chapter2.grothendieck-categories}{ }

% segment-id: o014.li2.chapter2.grothendieck-categories.p001
Semua pernyataan utama dalam bagian ini dirumuskan relatif terhadap pilihan
semesta Grothendieck, bukan dengan berbicara tentang kategori
``besar'' tanpa kualifikasi.

% segment-id: o014.li2.chapter2.grothendieck-categories.p002
Ingat pengertian limit kecil, kelengkapan, dan generator
(Definisi~\ref{def:generators}).

% segment-id: o014.li2.chapter2.grothendieck-categories.d001
\begin{definisi}[Kategori Grothendieck {\cite{Gr57}}]\label{def:Grothendieck-cat}
	\index{kategori Grothendieck (Grothendieck category)}
	Misalkan $\mathcal{A}$ kategori abelian. Kategori $\mathcal{A}$ disebut
	\emph{kategori Grothendieck} jika syarat-syarat berikut dipenuhi.
% segment-id: o014.li2.chapter2.grothendieck-categories.l001
	\begin{itemize}
% segment-id: o014.li2.chapter2.grothendieck-categories.i001
		\item $\mathcal{A}$ kolengkap; dengan kata lain, kategori ini memiliki
		semua $\varinjlim$ kecil;
% segment-id: o014.li2.chapter2.grothendieck-categories.i002
		\item $\mathcal{A}$ memiliki generator;
% segment-id: o014.li2.chapter2.grothendieck-categories.i003
		\item untuk setiap kategori kecil terfilter $I$
		(lihat \S\ref{sec:filtered-indlim}), funktor
		$\varinjlim:\mathcal{A}^I\to\mathcal{A}$ eksak. Secara ekuivalen,
		untuk setiap diagram $\alpha\to\beta\to\gamma$ dalam $\mathcal{A}^I$,
		berlaku
% segment-id: o014.li2.chapter2.grothendieck-categories.g001
		\begin{multline*}
			\forall i\in\Obj(I),\quad
			0\to\alpha(i)\to\beta(i)\to\gamma(i)\to0
			\quad\text{eksak} \\
			\implies 0\to\varinjlim\alpha\to\varinjlim\beta
			\to\varinjlim\gamma\to0\quad\text{eksak}.
		\end{multline*}
	\end{itemize}
\end{definisi}

% segment-id: o014.li2.chapter2.grothendieck-categories.p003
Contoh~\ref{eg:limit-exactness} telah menunjukkan bahwa $\varinjlim$
mempertahankan kokernel. Oleh karena itu, menurut
Catatan~\ref{rem:exactness-mono-epi}, syarat terakhir di atas juga ekuivalen
dengan pernyataan bahwa setiap $\varinjlim$ kecil terfilter mempertahankan
monomorfisme. Argumen berikut memperlihatkan kegunaan syarat ini.

% segment-id: o014.li2.chapter2.grothendieck-categories.pr001
\begin{proposisi}\label{prop:Grothendieck-cat-coprod-exact}
Misalkan $\mathcal{A}$ kategori Grothendieck. Untuk setiap himpunan kecil $I$,
pengambilan jumlah langsung (yakni koproduk) memberikan funktor eksak
$\bigoplus_I:\mathcal{A}^I\to\mathcal{A}$.
\end{proposisi}
% segment-id: o014.li2.chapter2.grothendieck-categories.pf001
\begin{bukti}
Kasus hingga jelas, sedangkan Proposisi~\ref{prop:limit-filtered-approx}
menyajikan $\bigoplus_I$ sebagai $\varinjlim$ terfilter dari jumlah-jumlah
langsung hingga.
\end{bukti}

% segment-id: o014.li2.chapter2.grothendieck-categories.p004
Selanjutnya kita memperoleh sifat yang tampak sederhana tetapi sangat berguna;
sifat ini juga bertumpu pada keeksakan $\varinjlim$ terfilter.

% segment-id: o014.li2.chapter2.grothendieck-categories.pr002
\begin{proposisi}\label{prop:Grothendieck-cat-intersection}
Tetapkan suatu himpunan terurut parsial kecil yang terfilter $(I,\leq)$.
Misalkan $\mathcal{A}$ kategori Grothendieck, atau secara lebih umum, misalkan
$\varinjlim:\mathcal{A}^{(I,\leq)}\to\mathcal{A}$ ada dan eksak. Untuk setiap
% segment-id: o014.li2.chapter2.grothendieck-categories.l002
\begin{compactitem}
% segment-id: o014.li2.chapter2.grothendieck-categories.i004
	\item $X\in\Obj(\mathcal{A})$ dan subobjek $Y\subset X$,
% segment-id: o014.li2.chapter2.grothendieck-categories.i005
	\item keluarga subobjek $\left(X_i\right)_{i\in I}$ dari $X$ yang memenuhi
	$i\leq j\implies X_i\subset X_j$,
\end{compactitem}
dengan notasi dari Konvensi~\ref{con:increasing-union}, kita mempunyai
% segment-id: o014.li2.chapter2.grothendieck-categories.g002
\begin{gather*}
	\varinjlim_{i\in I}X_i\rightiso\bigcup_{i\in I}X_i\subset X,\\
	Y\cap\left(\bigcup_{i\in I}X_i\right)
	=\bigcup_{i\in I}(Y\cap X_i)\;\in\mathrm{Sub}_X.
\end{gather*}
\end{proposisi}
% segment-id: o014.li2.chapter2.grothendieck-categories.pf002
\begin{bukti}
Karena $\varinjlim_i$ eksak, morfisme kanonik
$\iota:\varinjlim_iX_i\to X$ yang ditentukan oleh semua
$X_i\hookrightarrow X$ dan sifat universal $\varinjlim$ tetap merupakan
monomorfisme. Jika semua $X_i\hookrightarrow X$ berfaktor melalui suatu
subobjek $Z\subset X$, sifat universal tersebut memberikan faktorisasi $\iota$ sebagai
$\varinjlim_iX_i\to Z\subset X$. Jadi
$\iota:\varinjlim_iX_i\hookrightarrow X$ memang memberikan supremum dari
$(X_i)_{i\in I}$ di $\mathrm{Sub}_X$. Ini membuktikan persamaan pertama.

Selanjutnya, tuliskan morfisme kuosien $X\to X/Y$ sebagai $q$. Maka
$Y=\Ker(q)$ dan $Y\cap X_i=\Ker(q|_{X_i})$. Perhatikan keluarga morfisme yang
kompatibel $q|_{X_i}:X_i\to X/Y$. Karena pengambilan $\varinjlim_i$
mempertahankan kernel, bersama paragraf sebelumnya kita memperoleh
$Y\cap\left(\bigcup_{i\in I}X_i\right)=
\bigcup_{i\in I}(Y\cap X_i)$.
\end{bukti}

% segment-id: o014.li2.chapter2.grothendieck-categories.e001
\begin{contoh}[Kategori modul]
Misalkan $R$ gelanggang. Maka $R\dcate{Mod}$ adalah kategori Grothendieck.
Memang, sifat kolengkapnya sudah dikenal \cite[Teorema 6.2.2]{Li1}, sedangkan
keeksakan $\varinjlim$ terfilter diberikan oleh
\cite[Lema 6.8.3]{Li1}. Sebagai kasus khusus $R=\Z$, kategori $\cate{Ab}$
adalah kategori Grothendieck.
\end{contoh}

% segment-id: o014.li2.chapter2.grothendieck-categories.p005
Untuk mengembangkan teori ini lebih lanjut, kita memerlukan beberapa persiapan
mengenai generator.

% segment-id: o014.li2.chapter2.grothendieck-categories.d002
\begin{definisi}\index{generator!kuat (strong)}
Misalkan $s$ generator suatu kategori $\mathcal{C}$. Jika untuk setiap
monomorfisme $i:S_1\hookrightarrow S_2$ dalam $\mathcal{C}$, pemetaan terkait
$i_*:\Hom(s,S_1)\hookrightarrow\Hom(s,S_2)$ bijektif jika dan hanya jika $i$
isomorfisme, maka $s$ disebut \emph{generator kuat} dari $\mathcal{C}$.
\end{definisi}

% segment-id: o014.li2.chapter2.grothendieck-categories.pr003
\begin{proposisi}\label{prop:strong-generator-well-powered}
Misalkan $s$ merupakan generator kuat kategori $\mathcal{C}$ dan, untuk setiap
objek $X$ serta monomorfisme $S_i\hookrightarrow X$ dengan $i=1,2$, produk serat
$S_1\cap S_2:=S_1\dtimes{X}S_2$ ada. Jika
$\kappa:=|\Hom(s,X)|$, maka $|\mathrm{Sub}_X|\leq2^\kappa$.
\end{proposisi}
% segment-id: o014.li2.chapter2.grothendieck-categories.pf003
\begin{bukti}
Misalkan $s$ generator kuat dan $X$ sebarang objek. Akan ditunjukkan bahwa
% segment-id: o014.li2.chapter2.grothendieck-categories.g003
\begin{align*}
	\mathrm{Sub}_X&\longrightarrow
	\left\{\text{himpunan bagian dari }\Hom(s,X)\right\}\\
	S&\longmapsto\Hom(s,S)
\end{align*}
bersifat injektif, dan ini memberikan ketaksamaan yang diinginkan.

Ambil $S_1,S_2\in\mathrm{Sub}_X$ sedemikian sehingga
$\Hom(s,S_1)=\Hom(s,S_2)$. Jika $S_1\subset S_2$, definisi generator kuat
segera memberikan $S_1=S_2$. Dalam kasus umum,
$S_1\cap S_2\subset S_i$ untuk $i=1,2$ (lihat
Definisi~\ref{def:intersection-sum}), dan sebagai himpunan bagian dari
$\Hom(s,X)$ berlaku
% segment-id: o014.li2.chapter2.grothendieck-categories.q001
\[\Hom\left(s,S_1\cap S_2\right)
=\Hom(s,S_1)\cap\Hom(s,S_2)=\Hom(s,S_i),\qquad i=1,2.\]
Langkah sebelumnya memberikan $S_1=S_1\cap S_2=S_2$, sebagaimana diklaim.
\end{bukti}

% segment-id: o014.li2.chapter2.grothendieck-categories.pr004
\begin{proposisi}\label{prop:Abel-strong-generator}
Setiap generator dalam kategori abelian secara otomatis merupakan generator
kuat.
\end{proposisi}
% segment-id: o014.li2.chapter2.grothendieck-categories.pf004
\begin{bukti}
Misalkan $\mathcal{A}$ kategori abelian dan $s$ generatornya. Untuk
monomorfisme $i:S_1\hookrightarrow S_2$, perhatikan pasangan morfisme
% segment-id: o014.li2.chapter2.grothendieck-categories.g004
\[\begin{tikzcd}[column sep=large]
	S_2\arrow[shift left,r,"q:\text{ morfisme kuosien}"]
	\arrow[shift right,r,"0"']&S_2/S_1
\end{tikzcd}\]
Jika $i$ bukan isomorfisme, maka $q\neq0$. Jadi, dari definisi generator,
terdapat $\epsilon\in\Hom(s,S_2)$ sedemikian sehingga
$q\epsilon\neq0\epsilon=0$; morfisme $\epsilon$ ini tidak dapat berfaktor
melalui $S_1=\Ker(q)$.
\end{bukti}

% segment-id: o014.li2.chapter2.grothendieck-categories.c001
\begin{korolari}\label{prop:Grothendieck-cat-wellpowered}
Setiap kategori Grothendieck merupakan kategori berdaya-baik sekaligus
koberdaya-baik (Definisi~\ref{def:well-powered}).
\end{korolari}
% segment-id: o014.li2.chapter2.grothendieck-categories.pf005
\begin{bukti}
Kategori abelian memiliki semua produk serat hingga, sedangkan setiap
himpunan $\Hom$ kecil menurut asumsi. Jadi sifat berdaya-baik mengikuti dari
Proposisi~\ref{prop:strong-generator-well-powered} dan
Proposisi~\ref{prop:Abel-strong-generator}. Selain itu, dalam kategori abelian,
$\mathrm{Sub}_X$ dan $\mathrm{Quot}_X$ selalu memiliki kardinalitas yang sama.
\end{bukti}

% segment-id: o014.li2.chapter2.grothendieck-categories.c002
\begin{korolari}\label{prop:Grothendieck-cat-complete}
Setiap kategori Grothendieck $\mathcal{A}$ lengkap; dengan kata lain, kategori
ini memiliki semua $\varprojlim$ kecil.
\end{korolari}
% segment-id: o014.li2.chapter2.grothendieck-categories.pf006
\begin{bukti}
Kategori $\mathcal{A}$ kolengkap dan koberdaya-baik. Terapkan
Korolari~\ref{prop:SAFT-completeness}.
\end{bukti}

% segment-id: o014.li2.chapter2.grothendieck-categories.p006
Dengan demikian, kategori Grothendieck memiliki semua jumlah langsung kecil
dan produk langsung kecil. Latihan bab ini akan menunjukkan bahwa morfisme
kanonik $\delta:\bigoplus_{i\in I}X_i\to\prod_{i\in I}X_i$ dalam
\eqref{eqn:coprod-prod-delta} merupakan monomorfisme; untuk kategori modul,
hal ini tentu sudah jelas.

% segment-id: o014.li2.chapter2.grothendieck-categories.c003
\begin{korolari}\label{prop:Grothendieck-cat-rep1}
Misalkan $\mathcal{A}$ kategori Grothendieck. Funktor
$G:\mathcal{A}^{\opp}\to\cate{Set}$ terwakili jika dan hanya jika $G$
mempertahankan semua $\varprojlim$ kecil; secara lebih konkret, $G$ mengubah
setiap $\varinjlim$ kecil dalam $\mathcal{A}$ menjadi $\varprojlim$ kecil dalam
$\cate{Set}$.
\end{korolari}
% segment-id: o014.li2.chapter2.grothendieck-categories.pf007
\begin{bukti}
Terapkan Korolari~\ref{prop:cogenerator-representability} pada
$\mathcal{A}^{\opp}$.
\end{bukti}

% segment-id: o014.li2.chapter2.grothendieck-categories.p007
Perhatikan bahwa meskipun definisi kategori Grothendieck tidak swa-dual, versi
dual dari hasil di atas tetap berlaku; lihat
Korolari~\ref{prop:Grothendieck-cat-rep2}.

% segment-id: o014.li2.chapter2.grothendieck-categories.p008
Sekarang kita akan menunjukkan bahwa kategori Grothendieck memiliki cukup
banyak objek injektif. Bahkan, bukan hanya
monomorfisme $X\hookrightarrow I$ dalam Definisi~\ref{def:enough-injectives}
yang ada, melainkan juga bahwa $I$ dapat dipilih secara funktorial terhadap
$X$.

% segment-id: o014.li2.chapter2.grothendieck-categories.le001
\begin{lema}\label{prop:Grothendieck-injective-criterion}
Misalkan $\mathcal{A}$ kategori Grothendieck dan $s$ generatornya. Objek $I$
injektif jika dan hanya jika, untuk setiap monomorfisme $X\hookrightarrow s$,
setiap morfisme $X\to I$ dapat diperluas menjadi morfisme $s\to I$; dengan
kata lain, data tersebut dapat dilengkapi menjadi diagram komutatif
% segment-id: o014.li2.chapter2.grothendieck-categories.g005
$\begin{tikzcd}[row sep=small,column sep=small]
	X\arrow[r]\arrow[hookrightarrow,d]&I\\
	s\arrow[dashed,ru,"\exists"']&
\end{tikzcd}$.
\end{lema}
% segment-id: o014.li2.chapter2.grothendieck-categories.pf008
\begin{bukti}
Sifat injektif ekuivalen dengan pernyataan berikut: untuk setiap monomorfisme
$A\hookrightarrow B$, semua morfisme $f:A\to I$ dapat diperluas menjadi
morfisme $B\to I$. Jadi arah ``hanya jika'' jelas.

Sekarang kita buktikan arah sebaliknya. Ingat bahwa $\mathcal{A}$ adalah
kategori berdaya-baik. Untuk data $I\xleftarrow{f}A\hookrightarrow B$ seperti
di atas, perhatikan himpunan kecil
% segment-id: o014.li2.chapter2.grothendieck-categories.q002
\[\mathcal{S}:=\left\{\begin{array}{r|l}
	(A',f')&A\subset A'\subset B\quad\text{(sebagai subobjek)},\\
	&f':A'\to I\quad\text{memperluas }f
\end{array}\right\},\]
yang diurutkan parsial oleh relasi perluasan. Dari keeksakan $\varinjlim$
terfilter, setiap subhimpunan terurut total $\mathcal{T}$ dari $\mathcal{S}$
memiliki batas atas
% segment-id: o014.li2.chapter2.grothendieck-categories.q003
\[\widetilde{A}:=\varinjlim_{(A',f')\in\mathcal{T}}A',\qquad
\widetilde{f}:=\varinjlim_{(A',f')\in\mathcal{T}}f':\widetilde{A}\to I.\]
Lema Zorn menunjukkan bahwa $\mathcal{S}$ memiliki unsur maksimal $(A',f')$.
Kita hendak membuktikan bahwa $A'=B$. Jika tidak, karena $s$ generator kuat,
terdapat $g:s\to B$ yang tidak dapat berfaktor melalui $A'\hookrightarrow B$.
Kita akan menunjukkan bahwa $f'$ dapat diperluas menjadi
$A'+g(s)\to I$, bertentangan dengan kemaksimalan $(A',f')$.

Tetapkan $Y:=A'\cap g(s)$. Perhatikan bagian bergaris penuh dari diagram
komutatif berikut:
% segment-id: o014.li2.chapter2.grothendieck-categories.g006
\[\begin{tikzcd}[row sep=large]
	\Ker(g)\arrow[hookrightarrow,r]&g^{-1}(Y)
	\arrow[twoheadrightarrow,r,"g"]\arrow[hookrightarrow,d]
	&Y\arrow[hookrightarrow,r]\arrow[hookrightarrow,d]
	&A'\arrow[r,"{f'}"]&I\\
	&s\arrow[twoheadrightarrow,r,"g"']
	\arrow[dashed,rrru,crossing over,"\varphi" description]
	&g(s)\arrow[dashed,bend right,rru,"{f''}" description]&&
\end{tikzcd}\]
Menurut asumsi, terdapat $\varphi$ seperti ditunjukkan oleh garis putus-putus
yang membuat bagian segitiga komutatif. Karena $\varphi$ bernilai nol pada
$\Ker(g)$, terdapat pula $f''$ seperti pada garis putus-putus sehingga
$\varphi=f''g$. Dengan menarik balik sepanjang
$g^{-1}(Y)\twoheadrightarrow Y$ (yang merupakan epimorfisme berdasarkan
Proposisi~\ref{prop:Abel-cat-pull-push}), kita memperoleh bahwa pembatasan
$f'$ dan $f''$ pada $Y$ sama. Oleh sifat universal koproduk serat, keduanya
kemudian dapat direkatkan menjadi morfisme $A'+g(s)\to I$. Ini membuktikan
klaim.
\end{bukti}

% segment-id: o014.li2.chapter2.grothendieck-categories.p009
Sekarang kita akan memperkenalkan suatu varian dari apa yang disebut
``argumen objek kecil''. Pembaca sebaiknya meninjau sekilas
Definisi~\sourcecrossref{def:regular-cardinal}{A.2.7} dan pembahasan terkait tentang kardinal
reguler; uraian berikut memerlukan operasi serta sifat-sifat dasarnya.

% segment-id: o014.li2.chapter2.grothendieck-categories.d003
\begin{definisi}\label{def:I-small}
	\index{argumen objek kecil (small object argument)}
	Misalkan $\mathcal{C}$ kategori yang memiliki semua $\varinjlim$ kecil
	terfilter, $I\subset\Mor(\mathcal{C})$ suatu himpunan bagian, dan $\alpha$
	kardinal kecil reguler.
% segment-id: o014.li2.chapter2.grothendieck-categories.l003
	\begin{enumerate}[(i)]
% segment-id: o014.li2.chapter2.grothendieck-categories.i006
		\item Objek $X\in\Obj(\mathcal{C})$ disebut \emph{objek $\alpha$-kecil
		relatif terhadap $I$} jika syarat berikut dipenuhi. Misalkan
		$\widetilde{\alpha}$ kardinal kecil reguler dengan
		$\widetilde{\alpha}\geq\alpha$. Untuk setiap funktor
		$\beta\mapsto Y_\beta$ dari ordinal $\widetilde{\alpha}$, yang dipandang
		sebagai kategori kecil terfilter, ke $\mathcal{C}$, jika
		$Y_\beta\to Y_{\beta'}$ berada dalam $I$ untuk setiap
		$\beta\leq\beta'$, maka pemetaan kanonik
% segment-id: o014.li2.chapter2.grothendieck-categories.q004
		\[\varinjlim_\beta\Hom\left(X,Y_\beta\right)
		\longrightarrow\Hom\left(X,\varinjlim_\beta Y_\beta\right)\]
		(lihat Persamaan~\sourcecrossref{eqn:I-small-gen}{A.2.1}) merupakan bijeksi.

% segment-id: o014.li2.chapter2.grothendieck-categories.i007
		\item Jika $X$ merupakan objek $\alpha$-kecil relatif terhadap
		$\Mor(\mathcal{C})$, maka $X$ disebut \emph{objek $\alpha$-kecil}.
	\end{enumerate}
\end{definisi}

% segment-id: o014.li2.chapter2.grothendieck-categories.p010
Relatif terhadap $I$ yang telah dipilih, jika $\alpha\leq\alpha'$, maka sifat
$\alpha$-kecil mengakibatkan sifat $\alpha'$-kecil.

% segment-id: o014.li2.chapter2.grothendieck-categories.le002
\begin{lema}\label{prop:alpha-small}
Misalkan $\mathcal{A}$ kategori Grothendieck, $X$ suatu objeknya, dan $\alpha$
kardinal kecil reguler. Jika $\alpha>\kappa:=|\mathrm{Sub}_X|$, maka $X$
adalah objek $\alpha$-kecil relatif terhadap semua monomorfisme
(Definisi~\ref{def:I-small}).
\end{lema}
% segment-id: o014.li2.chapter2.grothendieck-categories.pf009
\begin{bukti}
Perhatikan data $\beta\mapsto Y_\beta$ dari Definisi~\ref{def:I-small}~(i),
dengan syarat bahwa $Y_\beta\to Y_{\beta'}$ merupakan monomorfisme apabila
$\beta\leq\beta'$. Tanpa mengurangi keumuman, kita dapat mengambil
$\alpha=\widetilde{\alpha}$ dalam definisi tersebut.

Karena $\varinjlim$ terfilter dalam $\mathcal{A}$ eksak, setiap
$Y_\beta\to\varinjlim_\gamma Y_\gamma$ tetap merupakan monomorfisme.
Keeksakan $\varinjlim$ terfilter dalam $\cate{Ab}$ juga menyiratkan bahwa
pemetaan
\[
\varinjlim_{\gamma<\alpha}\Hom(X,Y_\gamma)
\longrightarrow\Hom\left(X,\varinjlim_{\gamma<\alpha}Y_\gamma\right)
\]
injektif;
jadi tinggal membuktikan bahwa pemetaan itu surjektif. Selanjutnya, kita
pandang setiap $Y_\beta$ sebagai subobjek dari
$\varinjlim_\gamma Y_\gamma$.

Untuk suatu morfisme $f:X\to\varinjlim_\gamma Y_\gamma$, setiap
$\beta<\alpha$ menentukan subobjek $f^{-1}(Y_\beta)$ dari $X$. Kita dapat
mengasumsikan $f^{-1}(Y_\beta)\neq X$ untuk setiap $\beta$, sebab jika tidak,
tidak ada lagi yang perlu dibuktikan. Ingat bahwa $f^{-1}$ diberikan oleh
produk serat. Dengan sekali lagi menggunakan keeksakan $\varinjlim$
terfilter, kita memperoleh isomorfisme alami
% segment-id: o014.li2.chapter2.grothendieck-categories.q005
\[\varinjlim_\beta f^{-1}(Y_\beta)
\simeq f^{-1}\left(\varinjlim_\beta Y_\beta\right)=X.\]
Karena terdapat paling banyak $\kappa$ subobjek $f^{-1}(Y_\beta)$, ada suatu
himpunan bagian $S$ dari $\alpha$ dengan $|S|\leq\kappa$ sehingga ruas kiri
dapat diganti oleh $\varinjlim_{\beta\in S}$.

Perhatikan ordinal $\sigma:=\sup S\leq\alpha$. Akan ditunjukkan bahwa
$\sigma<\alpha$. Jika tidak, dari $|S|\leq\kappa<\alpha$ dan
$|S|\geq\mathrm{cf}(\alpha)$
(Proposisi~\sourcecrossref{prop:cf-generalities}{A.2.8}~(ii)) akan
diperoleh $\mathrm{cf}(\alpha)<\alpha$, bertentangan dengan kereguleran
$\alpha$. Klaim ini memastikan bahwa $f^{-1}(Y_\beta)$ selalu merupakan
subobjek dari $f^{-1}(Y_\sigma)$ untuk $\beta\in S$. Akibatnya,
$f^{-1}(Y_\sigma)=X$, sehingga $f$ berfaktor melalui $Y_\sigma$.
\end{bukti}

% segment-id: o014.li2.chapter2.grothendieck-categories.th001
\begin{teorema}[A.\ Grothendieck]\label{prop:Grothendieck-injective}
Misalkan $\mathcal{A}$ kategori Grothendieck. Tuliskan $\mathcal{I}$ untuk
subkategori penuh dari $\mathcal{A}$ yang terdiri atas objek-objek injektif,
dan tuliskan funktor inklusinya sebagai $\iota:\mathcal{I}\to\mathcal{A}$.
Terdapat funktor $F:\mathcal{A}\to\mathcal{I}$ beserta transformasi alami
$\varphi:\identity_{\mathcal{A}}\to\iota F$ sedemikian sehingga setiap
% segment-id: o014.li2.chapter2.grothendieck-categories.q006
\[\varphi_X:X\to F(X),\qquad X\in\Obj(\mathcal{A})\]
merupakan monomorfisme dalam $\mathcal{A}$. Khususnya, $\mathcal{A}$ memiliki
cukup banyak objek injektif.
\end{teorema}

% segment-id: o014.li2.chapter2.grothendieck-categories.pf010
\begin{bukti}
	Pilih suatu generator $s$. Tetapkan $F_0$ sebagai funktor
	$\identity_{\mathcal{A}}$. Langkah pertama ialah membentuk, untuk setiap
	objek $X$, diagram dorong keluar
% segment-id: o014.li2.chapter2.grothendieck-categories.g007
	\[\begin{tikzcd}
		\bigoplus_{Y \in \mathrm{Sub}_s} \bigoplus_{\varphi: Y \to X} Y \arrow[r] \arrow[hookrightarrow, d] & X \arrow[d, "{\varphi(0, 1)_X}"] \\
		\bigoplus_{Y \in \mathrm{Sub}_s} \bigoplus_{\varphi: Y \to X} s \arrow[r] & F_1(X) \arrow[phantom, lu, "\boxplus" description]
	\end{tikzcd}\]
	Proposisi \ref{prop:Abel-cat-pull-push} memastikan bahwa
	$\varphi(0,1)_X:X\to F_1(X)$ pada diagram tersebut merupakan
	monomorfisme. Ketika $X$ bervariasi, penetapan $X\mapsto F_1(X)$
	menjadi funktor dari $\mathcal{A}$ ke dirinya sendiri, dan
	$\varphi(0,1):F_0\to F_1$ menjadi transformasi alami.

	Secara lebih umum, untuk setiap ordinal kecil $\alpha$ kita akan membentuk
	funktor $F_\alpha:\mathcal{A}\to\mathcal{A}$ beserta keluarga transformasi
	$\varphi(\beta,\alpha):F_\beta\to F_\alpha$ untuk $\beta\leq\alpha$,
	dengan sifat-sifat berikut.
% segment-id: o014.li2.chapter2.grothendieck-categories.l004
	\begin{compactitem}
% segment-id: o014.li2.chapter2.grothendieck-categories.i008
		\item $\varphi(\alpha,\alpha)=\identity_{F_\alpha}$;
% segment-id: o014.li2.chapter2.grothendieck-categories.i009
		\item $\varphi(\beta,\alpha)_X:F_\beta(X)\to F_\alpha(X)$ merupakan
		monomorfisme untuk setiap $X$;
% segment-id: o014.li2.chapter2.grothendieck-categories.i010
		\item jika $\gamma\leq\beta\leq\alpha$, maka
		$\varphi(\beta,\alpha)\varphi(\gamma,\beta)=\varphi(\gamma,\alpha)$.
	\end{compactitem}
	Konstruksi ini dilakukan dengan rekursi transfinit
	\cite[\S 1.3]{Li1}. Mulai dari kasus $\alpha=0$, andaikan bahwa
	$F_\beta$ dan keluarga transformasi $\varphi(\beta',\beta)$ telah
	didefinisikan untuk setiap ordinal $\beta<\alpha$. Tetapkan
% segment-id: o014.li2.chapter2.grothendieck-categories.q007
	\[ F_\alpha(X) := \begin{cases}
		F_1(F_\beta(X)), & \alpha = \beta + 1:\ \text{ordinal penerus}, \\
		\varinjlim_{\beta < \alpha} F_\beta(X), & \alpha: \text{ordinal limit},
	\end{cases}\]
	dengan $\varinjlim$ dibentuk relatif terhadap
	$(\varphi(\beta',\beta))_{\beta'\leq\beta<\alpha}$. Sifat $F_1$ dan
	fakta bahwa $\varinjlim$ terfilter mempertahankan monomorfisme menentukan
	pilihan yang jelas bagi $\varphi(\beta,\alpha)$.

	Dengan Lema \sourcecrossref{prop:enough-regular-cardinal}{A.2.10}, pilih kardinal kecil reguler
	$\alpha$ sedemikian sehingga $\alpha>|\mathrm{Sub}_s|$. Kita tunjukkan
	bahwa $F_\alpha X$ injektif untuk setiap $X$. Diberikan
	$Y\hookrightarrow s$ dan $f:Y\to F_\alpha X$. Definisi kardinal reguler
	\sourcecrossref{def:regular-cardinal}{A.2.7} menyiratkan bahwa $\alpha$ merupakan ordinal
	limit. Dari konstruksi $F_\alpha$, inklusi
	$\mathrm{Sub}_Y\subset\mathrm{Sub}_s$, dan Lema
	\ref{prop:alpha-small}, diperoleh bahwa $f$ berfaktor melalui suatu
	$\varphi:Y\to F_\beta X$ dengan $\beta+1<\alpha$. Tinjau diagram
	komutatif
% segment-id: o014.li2.chapter2.grothendieck-categories.g008
	\[\begin{tikzcd}
		Y \arrow[r, "{\text{melalui}\; (Y, \varphi)}" inner sep=0.8em] \arrow[hookrightarrow, d] & \bigoplus_{Y' \in \mathrm{Sub}_s} \bigoplus_{\varphi': Y' \to F_\beta X} Y' \arrow[r] \arrow[hookrightarrow, d] & F_\beta X \arrow[d] \arrow[rd] & \\
		s \arrow[r, "{\text{melalui}\; (Y, \varphi)}"' inner sep=0.8em] & \bigoplus_{Y' \in \mathrm{Sub}_s} \bigoplus_{\varphi': Y' \to F_\beta X} s \arrow[r] & F_{\beta + 1} X \arrow[phantom, lu, "\boxplus" description] \arrow[r] & F_\alpha X
	\end{tikzcd}\]
	Dari diagram ini tampak langsung bahwa $f$ meluas menjadi morfisme
	$s\to F_\alpha X$. Lema
	\ref{prop:Grothendieck-injective-criterion} kemudian menyiratkan bahwa
	$F_\alpha X$ merupakan objek injektif. Akhirnya, ambil
	$F:=F_\alpha:\mathcal{A}\to\mathcal{I}$ dan
	$\varphi:=\varphi(0,\alpha)$.
\end{bukti}

% segment-id: o014.li2.chapter2.grothendieck-categories.c004
\begin{korolari}\label{prop:Grothendieck-cat-cogenerator}
	Setiap kategori Grothendieck mempunyai kogenerator injektif.
\end{korolari}
% segment-id: o014.li2.chapter2.grothendieck-categories.pf011
\begin{bukti}
	Pilih generator $s$ bagi kategori Grothendieck $\mathcal{A}$. Himpunan
	$\mathrm{Quot}_s$ diketahui kecil. Pilih objek injektif $I$ beserta
	monomorfisme
	$\bigoplus_{Q\in\mathrm{Quot}_s}Q\hookrightarrow I$. Kita gunakan
	kriteria Proposisi \ref{prop:injective-cogenerator} untuk menunjukkan bahwa
	$I$ merupakan kogenerator.

	Misalkan $T\in\Obj(\mathcal{A})$ tak nol. Ada morfisme tak nol
	$s\to T$, yang berfaktor sebagai $s\twoheadrightarrow Q'\hookrightarrow T$.
	Dengan demikian terdapat
% segment-id: o014.li2.chapter2.grothendieck-categories.q008
	\[ Q' \stackrel{\text{jelas}}{\hookrightarrow}
	\bigoplus_{Q \in \mathrm{Quot}_s} Q \hookrightarrow I. \]
	Menurut definisi objek injektif, komposisi ini meluas menjadi morfisme
	$T\to I$, dan morfisme tersebut jelas tak nol. Ini membuktikan klaim.
\end{bukti}

% segment-id: o014.li2.chapter2.grothendieck-categories.c005
\begin{korolari}\label{prop:Grothendieck-cat-rep2}
	Misalkan $\mathcal{A}$ kategori Grothendieck. Funktor
	$F:\mathcal{A}\to\cate{Set}$ terwakili jika dan hanya jika ia
	mempertahankan semua $\varprojlim$ kecil.
\end{korolari}
% segment-id: o014.li2.chapter2.grothendieck-categories.pf012
\begin{bukti}
	Kategori $\mathcal{A}$ mempunyai kogenerator; terapkan Korolari
	\ref{prop:cogenerator-representability}.
\end{bukti}

% segment-id: o014.li2.chapter2.grothendieck-categories.p011
Karakterisasi lain dari kategori Grothendieck diberikan oleh Teorema
Gabriel--Popescu--Kuhn \sourcecrossref{prop:GP}{A.3.2} dalam lampiran buku ini.

% Pemenggalan halaman dipaksakan di sini karena alasan tipografis.
\vfill
% segment-id: o014.li2.chapter2.grothendieck-categories.x001
\begin{Exercises}
% segment-id: o014.li2.chapter2.grothendieck-categories.ex001
	\item Tunjukkan bahwa subkategori aditif penuh dalam $\Z\dcate{Mod}$ yang
	dibentuk oleh modul-modul $\Z$ bebas bukan kategori abelian.

% segment-id: o014.li2.chapter2.grothendieck-categories.ex002
	\item Misalkan $X$ objek dalam kategori abelian dan
	$Y,Z\in\mathrm{Sub}_X$. Tinjau isomorfisme antarhimpunan terurut parsial
% segment-id: o014.li2.chapter2.grothendieck-categories.g009
	\[\begin{tikzcd}
		\left[Y \cap Z, Y \right] \arrow[d, "\simeq"] \arrow[r, "\sim"] & \left[ Z, Y+Z \right] \arrow[d, "\simeq"] \\
		\mathrm{Sub}_{Y/(Y \cap Z)} \arrow[r, "\sim"] & \mathrm{Sub}_{(Y+Z)/Z}
	\end{tikzcd}\]
	di mana anak panah vertikal berasal dari Teorema
	\ref{prop:Abel-cat-isom-thm} (ii), anak panah horizontal bawah berasal dari
	isomorfisme $Y/(Y\cap Z)\simeq(Y+Z)/Z$ pada Teorema
	\ref{prop:Abel-cat-isom-thm} (iii), dan anak panah horizontal atas adalah
	$W\mapsto W+Z$ dari Proposisi \ref{prop:lattice-diamond-isom}. Tunjukkan
	bahwa diagram tersebut komutatif.

% segment-id: o014.li2.chapter2.grothendieck-categories.ex003
	\item Misalkan $\Bbbk$ gelanggang komutatif dan $\mathcal{A}$ kategori
	abelian $\Bbbk$-linear. Buktikan bahwa jika $\Hom(X,X)$ merupakan modul
	Artinian atas $\Bbbk$ untuk setiap $X\in\Obj(\mathcal{A})$, maka
	$\mathcal{A}$ memenuhi syarat rantai ganda dalam Definisi
	\ref{def:bichain}.
% segment-id: o014.li2.chapter2.grothendieck-categories.hn001
	\begin{hint}
		Untuk rantai ganda $(X_n,\alpha_n,\beta_n)_{n=0}^\infty$, penetapan
		$f\mapsto\beta_n f\alpha_n$ menentukan barisan monomorfisme
		modul-$\Bbbk$
		$\Hom(X_{n+1},X_{n+1})\hookrightarrow\Hom(X_n,X_n)$. Tunjukkan bahwa
		morfisme ini merupakan isomorfisme jika dan hanya jika $\alpha_n$ dan
		$\beta_n$ keduanya merupakan isomorfisme.
	\end{hint}

	\index{kategori!Karoubi (Karoubian)}
% segment-id: o014.li2.chapter2.grothendieck-categories.ex004
	\item Misalkan $\mathcal{K}$ suatu kategori-$\cate{Ab}$ dengan objek nol.
	Jika $\Ker(e)$ ada untuk setiap $X\in\Obj(\mathcal{K})$ dan setiap
	idempoten $e\in\End(X)$, maka $\mathcal{K}$ disebut
	\emph{kategori Karoubi}.
% segment-id: o014.li2.chapter2.grothendieck-categories.l005
	\begin{enumerate}[(i)]
% segment-id: o014.li2.chapter2.grothendieck-categories.i011
		\item Buktikan bahwa setiap idempoten $e$ dalam kategori Karoubi
		mempunyai kokernel, serta terdapat isomorfisme kanonik
		$\Ker(e)\simeq\Coker(e)$.
% segment-id: o014.li2.chapter2.grothendieck-categories.i012
		\item Untuk setiap kategori-$\cate{Ab}$ $\mathcal{C}$ dengan objek nol,
		bentuk secara kanonik suatu kategori Karoubi $\mathrm{kar}(\mathcal{C})$
		beserta funktor aditif penuh dan setia
		$\varphi_{\mathcal{C}}:\mathcal{C}\to\mathrm{kar}(\mathcal{C})$
		sedemikian sehingga, untuk setiap kategori Karoubi $\mathcal{K}$, funktor
% segment-id: o014.li2.chapter2.grothendieck-categories.q009
		\[ \varphi_{\mathcal{C}}^*:
		\mathcal{K}^{\mathrm{kar}(\mathcal{C})}\to
		\mathcal{K}^{\mathcal{C}},\quad F\mapsto F\varphi_{\mathcal{C}} \]
		merupakan ekuivalensi. Di sini $\mathcal{K}^{\mathcal{C}}$ adalah
		kategori funktor yang terdiri dari semua funktor aditif
		$\mathcal{C}\to\mathcal{K}$, dan demikian pula untuk notasi lainnya.
% segment-id: o014.li2.chapter2.grothendieck-categories.i013
		\item Jika $\Bbbk$ gelanggang komutatif dan $\mathcal{C}$ bersifat
		$\Bbbk$-linear, maka $\mathrm{kar}(\mathcal{C})$ juga mempunyai struktur
		$\Bbbk$-linear yang alami.
	\end{enumerate}
% segment-id: o014.li2.chapter2.grothendieck-categories.p012
	Kategori $\mathrm{kar}(\mathcal{C})$ lazim disebut
	\emph{selubung Karoubi} dari $\mathcal{C}$; kategori ini diperoleh dengan
	menambahkan semua suku langsung pada $\mathcal{C}$ secara formal.

% segment-id: o014.li2.chapter2.grothendieck-categories.hn002
	\begin{hint}
		Ambil objek-objek $\mathrm{kar}(\mathcal{C})$ sebagai data $(X,e)$,
		dengan $X\in\Obj(\mathcal{C})$ dan $e\in\End(X)$ idempoten. Morfisme
		$(X,p)\to(Y,q)$ adalah morfisme $f:X\to Y$ dalam $\mathcal{C}$ yang
		memenuhi $qf=f=fp$, sedangkan
		$\varphi_{\mathcal{C}}(X)=(X,\identity_X)$.
	\end{hint}

% segment-id: o014.li2.chapter2.grothendieck-categories.ex005
	\item Untuk suatu gelanggang $R$, periksa bahwa semua modul-$R$ projektif
	membentuk kategori Karoubi yang bukan kategori abelian.

% segment-id: o014.li2.chapter2.grothendieck-categories.ex006
	\item Misalkan kategori abelian $\mathcal{A}$ mempunyai semua jumlah
	langsung kecil $\bigoplus_{i\in I}$ (atau semua produk kecil
	$\prod_{i\in I}$). Buktikan bahwa himpunan kecil
	$\Sigma\subset\Obj(\mathcal{A})$ merupakan keluarga generator (atau
	keluarga kogenerator) jika dan hanya jika, untuk setiap
	$X\in\Obj(\mathcal{A})$, terdapat jumlah langsung kecil (atau produk kecil)
		$S$ dari anggota-anggota $\Sigma$ beserta epimorfisme $S\twoheadrightarrow X$
	(atau monomorfisme $X\hookrightarrow S$).

% segment-id: o014.li2.chapter2.grothendieck-categories.hn003
	\begin{hint}
		Untuk keluarga generator, arah ``jika'' mudah. Untuk arah ``hanya jika'',
		ambil morfisme jelas
% segment-id: o014.li2.chapter2.grothendieck-categories.q010
		\[ f: \bigoplus_{s \in \Sigma} \bigoplus_{\varphi \in \Hom(s, X)} s \to X; \]
		setiap morfisme $\varphi$ berfaktor melalui $f$. Karena itu setiap
		komposisi $s\xrightarrow{\varphi'}X\to\Coker(f)$ bernilai $0$; simpulkan
		bahwa $\Coker(f)=0$.
	\end{hint}

% segment-id: o014.li2.chapter2.grothendieck-categories.ex007
	\item Misalkan $A,B,C$ subobjek dari objek $X$ dalam kategori abelian dan
	memenuhi $A\cap(B+C)=0=B\cap C$. Buktikan bahwa $A\cap B=0$ dan
	$(A+B)\cap C=0$.

% segment-id: o014.li2.chapter2.grothendieck-categories.hn004
	\begin{hint}
		Bentuk monomorfisme $(A+B)\cap C\to A$, lalu tunjukkan bahwa ia berfaktor
		melalui $A\cap(B+C)$. Anda dapat mencoba dahulu kasus kategori modul.
	\end{hint}

% segment-id: o014.li2.chapter2.grothendieck-categories.ex008
	\item Buktikan bahwa kategori abelian modul-$\Z$ yang dibangkitkan secara hingga
	mempunyai cukup banyak objek projektif, tetapi tidak mempunyai cukup banyak
	objek injektif.

% segment-id: o014.li2.chapter2.grothendieck-categories.ex009
	\item (Lema Schanuel) Tinjau objek projektif $P,Q$ dalam suatu kategori
	abelian dan barisan eksak pendek
% segment-id: o014.li2.chapter2.grothendieck-categories.g010
	\begin{gather*}
		0 \to K \to P \xrightarrow{\phi} M \to 0, \\
		0 \to L \to Q \xrightarrow{\psi} M \to 0.
	\end{gather*}
	Dari data ini, bentuk diagram tarik balik
% segment-id: o014.li2.chapter2.grothendieck-categories.g011
	\[\begin{tikzcd}
		X \arrow[r, "{\psi'}"] \arrow[d, "{\phi'}"'] & P \arrow[d, "\phi"] \\
		Q \arrow[r, "\psi"'] & M \arrow[phantom, lu, "\Box" description] .
	\end{tikzcd}\]
	Tunjukkan bahwa terdapat barisan eksak pendek
	$0\to L\to X\xrightarrow{\psi'}P\to0$ dan
	$0\to K\to X\xrightarrow{\phi'}Q\to0$, lalu buktikan bahwa terdapat
	isomorfisme $K\oplus Q\simeq L\oplus P$.
% segment-id: o014.li2.chapter2.grothendieck-categories.hn005
	\begin{hint} Terapkan Proposisi \ref{prop:pull-back-ker} dan
	\ref{prop:Abel-cat-pull-push}. \end{hint}

% segment-id: o014.li2.chapter2.grothendieck-categories.ex010
	\item Buktikan bahwa $\Q/\Z$ merupakan kogenerator injektif bagi
	$\cate{Ab}$.

% segment-id: o014.li2.chapter2.grothendieck-categories.ex011
	\item Periksa secara rinci relasi adjoin dalam Contoh
	\ref{eg:A2-injective-projective}.

% segment-id: o014.li2.chapter2.grothendieck-categories.ex012
	\item Misalkan $\mathcal{C}$ kategori tak kosong dan $\mathcal{A}$ kategori
	abelian. Untuk setiap $c\in\Obj(\mathcal{C})$, funktor evaluasi
	$\mathrm{ev}_c:\mathcal{A}^{\mathcal{C}}\to\mathcal{A}$ memetakan $F$
	ke $Fc$.
% segment-id: o014.li2.chapter2.grothendieck-categories.l006
	\begin{enumerate}[(i)]
% segment-id: o014.li2.chapter2.grothendieck-categories.i014
		\item Andaikan $\mathcal{A}$ mempunyai semua jumlah langsung berbentuk
		$\bigoplus_{c\in\Obj(\mathcal{C})}$ dan
		$\bigoplus_{f\in\Hom_{\mathcal{C}}(c,c')}$. Untuk setiap
		$c\in\Obj(\mathcal{C})$, definisikan funktor
		$\mathcal{L}_c:\mathcal{A}\to\mathcal{A}^{\mathcal{C}}$ pada objek
		dengan
		$(\mathcal{L}_cX)(c')=\bigoplus_{\Hom(c,c')}X$. Lengkapi definisinya pada
		morfisme sehingga $\mathcal{L}_c$ menjadi adjoin kiri bagi
		$\mathrm{ev}_c$.
%		\begin{hint}
%			Diberikan $c'\to c''$, definisi funktor pada morfisme ditentukan oleh
%			$\Hom(c,c')\to\Hom(c,c'')$ dan $\identity_X$.
%		\end{hint}
% segment-id: o014.li2.chapter2.grothendieck-categories.i015
		\item Dengan asumsi di atas, buktikan bahwa jika $\mathcal{A}$ mempunyai
		cukup banyak objek projektif, maka $\mathcal{A}^{\mathcal{C}}$ juga
		demikian.
% segment-id: o014.li2.chapter2.grothendieck-categories.i016
		\item Selidiki versi dualnya. Andaikan $\mathcal{A}$ mempunyai semua
		produk berbentuk $\prod_{c\in\Obj(\mathcal{C})}$ dan
		$\prod_{f\in\Hom_{\mathcal{C}}(c',c)}$. Definisikan funktor
		$\mathcal{R}_c:\mathcal{A}\to\mathcal{A}^{\mathcal{C}}$ dengan
		$(\mathcal{R}_cX)(c')=\prod_{\Hom(c',c)}X$ yang menjadi adjoin kanan
		bagi $\mathrm{ev}_c$. Selanjutnya buktikan bahwa jika $\mathcal{A}$
		mempunyai cukup banyak objek injektif, maka
		$\mathcal{A}^{\mathcal{C}}$ juga demikian.
% segment-id: o014.li2.chapter2.grothendieck-categories.i017
		\item Selidiki hubungan konstruksi-konstruksi tersebut dengan Contoh
		\ref{eg:A2-injective-projective} ketika $\mathcal{C}=\mathbf{2}$.
	\end{enumerate}

% segment-id: o014.li2.chapter2.grothendieck-categories.ex013
	\item Dengan notasi soal sebelumnya, misalkan $\mathcal{C}$ kategori kecil.
	Buktikan bahwa jika $s$ merupakan generator (atau kogenerator) bagi
	$\mathcal{A}$, maka semua $\mathcal{L}_c(s)$ (atau $\mathcal{R}_c(s)$),
	dengan $c$ merentang $\Obj(\mathcal{C})$, membentuk keluarga generator
	(atau keluarga kogenerator) bagi $\mathcal{A}^{\mathcal{C}}$, asalkan
	jumlah langsung (atau produk langsung) yang diperlukan ada.

% segment-id: o014.li2.chapter2.grothendieck-categories.hn006
	\begin{hint}
		Untuk kasus generator, bagi setiap $c\in\Obj(\mathcal{C})$ dan setiap
		morfisme $f:X\to Y$ dalam $\mathcal{A}^{\mathcal{C}}$ terdapat diagram
		komutatif
% segment-id: o014.li2.chapter2.grothendieck-categories.g012
		\[\begin{tikzcd}
			\Hom_{\mathcal{A}}(s, Y(c)) \arrow[r, "\sim"] & \Hom_{\mathcal{A}^{\mathcal{C}}}(\mathcal{L}_c(s), Y) \\
			\Hom_{\mathcal{A}}(s, X(c)) \arrow[r, "\sim"'] \arrow[u, "{f(c)_*}"] & \Hom_{\mathcal{A}^{\mathcal{C}}}(\mathcal{L}_c(s), X) \arrow[u, "{f_*}"'] .
		\end{tikzcd}\]
		Syarat agar suatu keluarga menjadi keluarga generator setara dengan
		pernyataan bahwa, ketika
		$c$ merentang $\Obj(\mathcal{C})$, semua $f_*$ pada kolom kedua
		menentukan $f$ secara tunggal.
	\end{hint}

% segment-id: o014.li2.chapter2.grothendieck-categories.ex014
	\item (Konstruksi langsung kuosien Serre) Untuk subkategori Serre
	$\mathcal{T}$ dari kategori abelian $\mathcal{A}$, definisikan kategori baru
	yang himpunan objeknya sama dengan $\Obj(\mathcal{A})$, sedangkan himpunan
	morfisme dari $X$ ke $Y$ ialah
% segment-id: o014.li2.chapter2.grothendieck-categories.q011
	\[ \varinjlim_{X' \subset X, Y' \subset Y} \Hom_{\mathcal{A}}(X', Y/Y'),
	\quad \text{dengan}\; X/X',\;Y'\in\Obj(\mathcal{T}). \]
% segment-id: o014.li2.chapter2.grothendieck-categories.l007
	\begin{enumerate}[(i)]
% segment-id: o014.li2.chapter2.grothendieck-categories.i018
		\item Lengkapi definisi komposisi morfisme dan tunjukkan bahwa kategori ini
		ekuivalen dengan kuosien Serre $\mathcal{A}/\mathcal{T}$ dalam Teorema
		\ref{prop:Serre-quotient}.
% segment-id: o014.li2.chapter2.grothendieck-categories.hn007
		\begin{hint}
			% Rujukan: Gelfand--Manin, hlm. 121.
			Periksa bahwa kategori ini mempunyai sifat universal yang
			dinyatakan dalam Teorema \ref{prop:Serre-quotient}.
		\end{hint}
% segment-id: o014.li2.chapter2.grothendieck-categories.i019
		\item Gunakan konstruksi ini untuk menunjukkan bahwa jika $\mathcal{A}$
		merupakan kategori berdaya-baik (Definisi \ref{def:well-powered}), maka
		$\mathcal{A}/\mathcal{T}$ juga merupakan kategori-$\mathcal{U}$, dengan
		$\mathcal{U}$ semesta Grothendieck yang telah dipilih.
	\end{enumerate}
	\index{masalah ukuran (size issue)}

% segment-id: o014.li2.chapter2.grothendieck-categories.ex015
	\item Misalkan $\mathcal{B}$ subkategori abelian dari $\mathcal{A}$, dan
	untuk setiap $X\in\Obj(\mathcal{A})$ terdapat barisan subobjek
	$0=X_0\subset\cdots\subset X_n=X$ sedemikian sehingga
	$X_i/X_{i-1}\in\Obj(\mathcal{B})$. Buktikan bahwa funktor inklusi
	$\mathcal{B}\to\mathcal{A}$ menginduksi isomorfisme grup
	$\mathrm{K}_0(\mathcal{B})\rightiso\mathrm{K}_0(\mathcal{A})$.

% segment-id: o014.li2.chapter2.grothendieck-categories.hn008
	\begin{hint}
		Gunakan Teorema Penghalusan Schreier
		\ref{prop:Schreier-refinement} untuk menunjukkan bahwa
		$[X]\mapsto\sum_{i=1}^n[X_i/X_{i-1}]$ memberikan pemetaan invers yang
		tidak bergantung pada pilihan barisan subobjek.
	\end{hint}

% segment-id: o014.li2.chapter2.grothendieck-categories.ex016
	\item Misalkan $\mathrm{isom}$ adalah himpunan kelas isomorfisme semua objek
	dalam kategori aditif $\mathcal{A}$. Tuliskan $\lrangle{X}$ untuk citra kelas
	isomorfisme objek $X$ dalam modul $\Z$ bebas
	$\Z^{\oplus\mathrm{isom}}$. Definisikan versi jumlah langsung dari grup
	$\mathrm{K}_0$ dengan
% segment-id: o014.li2.chapter2.grothendieck-categories.q012
	\[ \mathrm{K}_{\oplus}(\mathcal{A}) := \Z^{\oplus \mathrm{isom}} \big/
	\text{submodul yang dibangun oleh semua}\;
	\lrangle{X \oplus Y} - \lrangle{X} - \lrangle{Y}. \]

	Misalkan pula $\mathrm{ind}$ adalah himpunan kelas isomorfisme objek tak
	terurai. Buktikan bahwa jika $\mathcal{A}$ kategori Karoubi dan setiap objek
	tak nol mempunyai dekomposisi jumlah langsung hingga
	$X\simeq X_1\oplus\cdots\oplus X_k$, dengan $X_i\neq0$ dan
	$\End(X_i)$ gelanggang lokal (lihat pembahasan sebelum Korolari
	\ref{prop:indecomposable-local-ring}), maka terdapat isomorfisme kanonik
% segment-id: o014.li2.chapter2.grothendieck-categories.q013
	\[ \Z^{\oplus \mathrm{ind}} \rightiso \mathrm{K}_{\oplus}(\mathcal{A}). \]

% segment-id: o014.li2.chapter2.grothendieck-categories.hn009
	\begin{hint}
		Ingat bahwa jika $\End(X_i)$ gelanggang lokal, maka $X_i$
		tak terurai; buktikan langsung atau lihat \cite[Lema 6.12.6]{Li1}.
		Selanjutnya gunakan \cite[Teorema 6.12.8]{Li1} untuk memperoleh keunikan
		dekomposisi.
	\end{hint}

% segment-id: o014.li2.chapter2.grothendieck-categories.ex017
	\item Misalkan $\mathcal{A}$ kategori Grothendieck dan
	$(X_i)_{i\in I}$ suatu keluarga objek dengan $I$ himpunan kecil. Buktikan
	bahwa morfisme kanonik
	$\delta:\bigoplus_{i\in I}X_i\to\prod_{i\in I}X_i$ dalam
	\eqref{eqn:coprod-prod-delta} merupakan monomorfisme.

% segment-id: o014.li2.chapter2.grothendieck-categories.hn010
	\begin{hint}
		Untuk setiap himpunan bagian hingga $F\subset I$, tinjau diagram komutatif
% segment-id: o014.li2.chapter2.grothendieck-categories.g013
		\[\begin{tikzcd}
			\bigoplus_{i \in I} X_i \arrow[r, "\delta"] & \prod_{i \in I} X_i \\
			\bigoplus_{i \in F} X_i \arrow[r, "\sim", "\delta_F"'] \arrow[hookrightarrow, u] & \prod_{i \in F} X_i \arrow[hookrightarrow, u, "{\iota_F}"']
		\end{tikzcd}\]
		Ambil $\varinjlim$ terfilter dari
		$\iota_F\delta_F$ untuk semua $F\subset I$ hingga, lalu simpulkan bahwa
		$\delta$ monomorfisme; lihat Proposisi
		\ref{prop:limit-filtered-approx}.
	\end{hint}

% segment-id: o014.li2.chapter2.grothendieck-categories.ex018
	\item Buktikan bahwa jika $\mathcal{A}$ kategori Grothendieck, maka kategori
	kompleks $\cate{C}(\mathcal{A})$ (lihat
	\S\sourcecrossref{sec:additive-cplx}{3.1}) dan kategori funktor
	$\mathcal{A}^{\mathcal{C}}$ keduanya merupakan kategori Grothendieck, dengan
	$\mathcal{C}$ sebarang kategori kecil.

% segment-id: o014.li2.chapter2.grothendieck-categories.ex019
	\item (Karakterisasi jumlah langsung) Misalkan $\mathcal{A}$ kategori
	Grothendieck. Tinjau $X\in\Obj(\mathcal{A})$ beserta keluarga subobjek
	$(X_i)_{i\in I}$, dengan $I$ himpunan kecil. Morfisme kanonik
	$\bigoplus_{i\in I}X_i\to X$ memiliki pembatasan pada setiap $X_i$ yang
	merupakan inklusi
	$X_i\hookrightarrow X$. Buktikan bahwa
	$\bigoplus_{i\in I}X_i\rightiso X$ jika dan hanya jika sifat-sifat berikut
	terpenuhi.
% segment-id: o014.li2.chapter2.grothendieck-categories.l008
	\begin{enumerate}[(i)]
% segment-id: o014.li2.chapter2.grothendieck-categories.i020
		\item $X=\sum_{i\in I}X_i$;
% segment-id: o014.li2.chapter2.grothendieck-categories.i021
		\item untuk setiap $i\in I$ berlaku
		$X_i\cap\sum_{\substack{j\in I\\j\neq i}}X_j=0$.
	\end{enumerate}
	Perhatikan bahwa syarat (ii) cukup diperiksa pada himpunan bagian hingga
	dari $I$. Jika $I$ sendiri hingga, pernyataan di atas berlaku dalam sebarang
	kategori abelian.

% segment-id: o014.li2.chapter2.grothendieck-categories.hn011
	\begin{hint}
		Sifat (i) setara dengan surjektivitas morfisme kanonik, sedangkan (ii)
		setara dengan injektivitasnya. Yang pertama langsung; untuk yang kedua,
		gunakan $\varinjlim$ terfilter untuk mereduksi ke kasus $I$ hingga.
	\end{hint}

% segment-id: o014.li2.chapter2.grothendieck-categories.ex020
	\item Tinjau objek $X$ dalam kategori Grothendieck $\mathcal{A}$.
% segment-id: o014.li2.chapter2.grothendieck-categories.l009
	\begin{enumerate}[(i)]
% segment-id: o014.li2.chapter2.grothendieck-categories.i022
		\item Lengkapi bukti terperinci dari Catatan
		\ref{rem:Grothendieck-ss-decomp}.
% segment-id: o014.li2.chapter2.grothendieck-categories.hn012
		\begin{hint}
			Gunakan latihan sebelumnya untuk menangani jumlah langsung tak hingga.
		\end{hint}
% segment-id: o014.li2.chapter2.grothendieck-categories.i023
		\item Dengan asumsi berikut, buktikan bahwa jika $X$ terbelah maka $X$
		semisederhana: untuk setiap $Y\in\Obj(\mathcal{A})$ tak nol, terdapat
		subobjek tak nol $Y_0\subset Y$ sedemikian sehingga setiap rantai dalam
		himpunan terurut parsial
		$\mathrm{Sub}_{Y_0}\smallsetminus\{Y_0\}$ mempunyai batas atas.
		Periksa asumsi ini untuk $\mathcal{A}=R\dcate{Mod}$.
% segment-id: o014.li2.chapter2.grothendieck-categories.hn013
		\begin{hint}
			Anda dapat mengikuti bagian (iii) $\implies$ (i) dalam bukti
			\cite[Proposisi 6.11.4]{Li1}, serta bukti
			\cite[Lema 6.11.3]{Li1}.
		\end{hint}
	\end{enumerate}
\end{Exercises}
